\chapter{Min-max trees and the $\PsiOp{i}$ operators}
\label{ch:mmtree}

This chapter formalizes Stanley EC1 \S1.6.3, pp.~56--57: the min-max
tree construction $\Mtree{w}$, its window structure, and the
involutive operators $\PsiOp{i}$ that generate the
\emph{$M$-equivalence class} $\Mclass{w}$ on which the cd-monomial
$\PhiW{w}$ will be defined (Chapter~\ref{ch:cdindex}).

% ============================================================
\section{The min-max tree of a sequence}
% ============================================================

\begin{definition}[Min-max tree datatype]
  \label{def:mmtree}
  \rocq{mathcomp_eulerian.mmtree.mmtree}
  \rocqok
  A \emph{min-max tree} over a type $T$ is the inductive
  \[
    \mathrm{mmtree}\,T \;::=\;
      \mathrm{Leaf}
      \;\mid\;
      \mathrm{Node}\,(\mathrm{mmtree}\,T)\,T\,(\mathrm{mmtree}\,T).
  \]
  Its in-order traversal is
  \[
    \mathrm{mmtree\_to\_seq}(\mathrm{Node}\,\ell\,x\,r)
      \;=\;
      \mathrm{mmtree\_to\_seq}(\ell) \mathbin{+\!\!+} x :: \mathrm{mmtree\_to\_seq}(r).
  \]
\end{definition}

\begin{definition}[The mm position]
  \label{def:mm_pos}
  \rocq{mathcomp_eulerian.psi_core.mm_pos}
  \rocqok
  For a non-empty sequence $s$, $\code{mm\_pos}(s)$ is the position of
  $s$'s minimum if it occurs strictly before its maximum, and the
  position of its maximum otherwise.  This is the position of the root
  of $\Mtree{s}$.
\end{definition}

\begin{lemma}[mm position is in range]
  \label{lem:mm_pos_lt}
  \uses{def:mm_pos}
  \rocq{mathcomp_eulerian.psi_core.mm_pos_lt}
  \rocqok
  For non-empty $s$, $\code{mm\_pos}(s) < |s|$.
\end{lemma}

\begin{theorem}[Min-max tree round-trip]
  \label{thm:mmtree_of_seqK}
  \uses{def:mmtree, def:mm_pos}
  \rocq{mathcomp_eulerian.mmtree.mmtree_of_seqK}
  \rocqok
  For every $s \in \mathrm{seq}\,\NN$,
  \[
    \mathrm{mmtree\_to\_seq}(\mathrm{mmtree\_of\_seq}(s)) = s.
  \]
  This identifies $\Mtree{s}$ as a faithful representation of the
  underlying sequence: every position in $w$ corresponds to a $\mathrm{Node}$
  of $\Mtree{w}$.
\end{theorem}

% ============================================================
\section{Windows, internal vertices, and left children}
% ============================================================

For a position $i$ in $w$, its \emph{window} is the interval
$[i,\; i + \windowSize_i(w))$, the in-order range of the subtree
rooted at the $\mathrm{Node}$ for position $i$ in $\Mtree{w}$.

\begin{definition}[Window size]
  \label{def:window_size}
  \uses{def:mm_pos}
  \rocq{mathcomp_eulerian.psi_core.window_size}
  \rocqok
  $\windowSize_i(w)$ is defined by fuel-based recursion that descends
  $\Mtree{w}$ via $\code{mm\_pos}$: at the root position one returns
  $|w| - i$; otherwise one recurses into the left or right subtree
  according to whether $i$ lies before or after the root.
\end{definition}

\begin{definition}[Internal vertex]
  \label{def:is_internal}
  \uses{def:window_size}
  \rocq{mathcomp_eulerian.psi_cdindex_defs.is_internal}
  \rocqok
  Position $i$ is \emph{internal} in $w$ iff $i < |w|$ and
  $\windowSize_i(w) > 1$.  Equivalently, the $\mathrm{Node}$ for
  position $i$ in $\Mtree{w}$ has at least one non-$\mathrm{Leaf}$
  child.
\end{definition}

\begin{definition}[Has left child]
  \label{def:has_left_child}
  \uses{def:is_internal, def:window_size}
  \rocq{mathcomp_eulerian.psi_descent_v2.has_left_child}
  \rocqok
  $\hasLC_i(w)$ is true iff in $\Mtree{w}$ the $\mathrm{Node}$ at
  position $i$ has a non-$\mathrm{Leaf}$ left subtree.  This is
  Stanley's $\hasLC$ predicate from p.~57 (it distinguishes the
  cd-letters \texttt{c} and \texttt{d}; see Chapter~\ref{ch:cdindex}).
\end{definition}

\begin{definition}[Internal vertices of $w$]
  \label{def:internal_vertices}
  \uses{def:is_internal}
  \rocq{mathcomp_eulerian.psi_cdindex_defs.internal_vertices}
  \rocqok
  $\code{internal\_vertices}(w)$ is the strictly increasing sequence
  of positions $i$ for which $\internal_i(w)$ holds.
\end{definition}

% ============================================================
\section{The operators $\PsiOp{i}$}
% ============================================================

For each internal position $i$, $\PsiOp{i}$ is a permutation on
sequences that toggles the \emph{local order} of the window at $i$
without changing the unlabelled tree shape.  This is Stanley's
$\psi_i$ from EC1 \S1.6.3, p.~57.

\begin{definition}[The $\PsiOp{i}$ operator]
  \label{def:psi}
  \uses{def:window_size, def:has_left_child, def:internal_vertices}
  \rocq{mathcomp_eulerian.psi_core.psi}
  \rocqok
  For an internal position $i$ of $w$, $\PsiOp{i}(w)$ is defined by a
  rank-shift inside the window at $i$ (the precise formula is given
  in \code{psi\_core.v:219}).  For non-internal $i$, $\PsiOp{i}$ acts
  as the identity.
\end{definition}

\begin{theorem}[Each $\PsiOp{i}$ is an involution]
  \label{thm:psi_involutive}
  \uses{def:psi}
  \rocq{mathcomp_eulerian.psi_core.psi_involutive}
  \rocqok
  For every $i \in \NN$ and every uniq sequence $w$,
  $\PsiOp{i}(\PsiOp{i}(w)) = w$.
\end{theorem}

\begin{theorem}[Stanley Fact \#1: commutativity of $\PsiOp{i}$]
  \label{thm:psi_comm}
  \uses{def:psi, thm:psi_involutive}
  \rocq{mathcomp_eulerian.psi_comm.psi_comm}
  \rocqok
  For all $i, j \in \NN$ and every uniq sequence $w$,
  $\PsiOp{i}(\PsiOp{j}(w)) = \PsiOp{j}(\PsiOp{i}(w))$.
  This is the commutativity half of Stanley's
  \emph{Fact \#1} (EC1 p.~57): the $\PsiOp{i}$ for internal vertices
  generate an abelian group $G_w \cong (\ZZ/2)^{\iota(w)}$, where
  $\iota(w)$ is the number of internal vertices.
\end{theorem}

\begin{definition}[Iterated $\Psi$ application]
  \label{def:apply_psis}
  \uses{def:psi}
  \rocq{mathcomp_eulerian.psi_cdindex_defs.apply_psis}
  \rocqok
  For a sequence $\mathrm{ss}$ of internal positions,
  $\code{apply\_psis}(\mathrm{ss})(w)$ is the left-fold composition of
  the corresponding $\PsiOp{i}$.  By
  Theorem~\ref{thm:psi_comm} and Theorem~\ref{thm:psi_involutive},
  the result depends only on the \emph{set} $\mathrm{ss} \subseteq
  \code{internal\_vertices}(w)$, and the M-equivalence class
  $\Mclass{w} = \{\code{apply\_psis}(\mathrm{ss})(w) : \mathrm{ss}
  \subseteq \code{internal\_vertices}(w)\}$ has size $2^{\iota(w)}$.
\end{definition}

\begin{lemma}[$\Psi$ preserves the unlabelled tree shape]
  \label{lem:shape_apply_psis}
  \uses{def:apply_psis, def:window_size, def:has_left_child, def:is_internal, def:internal_vertices}
  \rocq{mathcomp_eulerian.psi_cdindex_core.window_size_apply_psis}
  \rocq{mathcomp_eulerian.psi_cdindex_core.has_left_child_apply_psis}
  \rocq{mathcomp_eulerian.psi_cdindex_core.is_internal_apply_psis}
  \rocq{mathcomp_eulerian.psi_cdindex_core.internal_vertices_apply_psis}
  \rocqok
  For every $\mathrm{ss}$, every position $i$, and every uniq $w$:
  $\windowSize_i(\code{apply\_psis}(\mathrm{ss})(w)) = \windowSize_i(w)$,
  $\hasLC_i(\code{apply\_psis}(\mathrm{ss})(w)) = \hasLC_i(w)$,
  $\internal_i(\code{apply\_psis}(\mathrm{ss})(w)) = \internal_i(w)$,
  and the lists of internal vertices coincide.

  These four invariances express that $\Psi$-action permutes labels
  inside windows but preserves the unlabelled tree shape — the
  load-bearing geometric fact behind Stanley Fact~\#1's
  well-definedness clause for $\PhiW{w}$ (Chapter~\ref{ch:cdindex}).
\end{lemma}
